\sectiontitle{一、单项选择题}{10}{20}

% A01 S07:10 | 基础
\mc{ETPP 使用 $f_V=g_V+[1+H_{IE}(V_c,V_n)]D_V(V_c,V_n)$，其中 $H_{IE}(V_c,V_n)=\min[\chi_1e^{-D_{free}(V_c)/c_V},\chi_2e^{-D_{free}(V_n)/c_V}]$。已知 $g_V=10$、$D_V=5$、$c_V=1$、$D_{free}(V_c)=0$、$D_{free}(V_n)=\ln 2$、$\chi_1=\chi_2=1$。计算 $f_V$。}
{$12.5$}{$17.5$}{$20$}{$15$}
{B}{先算 $H_{IE}=\min(1,e^{-\ln2})=0.5$。因此 $f_V=10+(1+0.5)\times5=17.5$。}

% A02 S07:11 | 基础
\mc{一个地图有 4 个单元，未知单元的概率为 $(0.5,0.5,0,0)$。按 $H=-\sum_i p_i\ln p_i$、$H_{max}=\ln4$ 和 $I=H/H_{max}$ 计算归一化信息熵。停止条件为 $I\le I_{thr}$，且 $I_{thr}=0.6$。本轮应如何处理？}
{继续探索，因为熵大于阈值}{先把 $H$ 与 $I_{thr}$ 比较}{停止探索，因为 $I=0.5$}{不能判断，因为概率中有零}
{C}{$H=-2(0.5\ln0.5)=\ln2$，$H_{max}=\ln4$，故 $I=\ln2/\ln4=0.5\le0.6$。零概率项按 $0\ln0=0$ 处理。}

% A03 S07:10 | 基础
\mc{ACXCOOT 的自适应迭代次数按 $Iternum=b_2n_{pf}/(b_1e)$ 计算，并以 $Iter=\min(Iternum,Iter_{max})$ 限幅。给定 $b_1=2$、$b_2=3$、$n_{pf}=40$、$Iter_{max}=20$，实际迭代次数应取哪一个值？}
{$22.07$}{$2.21$}{$40$}{$20$}
{D}{未限幅值为 $3\times40/(2e)\approx22.07$。它超过 $Iter_{max}$，所以实际值取 $20$。}

% A04 S07:20 | 基础
\mc{稀疏注意力行的 logits 为 $(1,3,0)$，掩码为 $(1,1,0)$。仅对掩码为 1 的项应用 $\operatorname{sparse\_softmax}(z_i)=e^{z_i-1}/\sum_j e^{z_j-1}$。输出应为哪一项？}
{$(0.5,0.5,0)$}{$(0.119,0.881,0)$}{$(0.119,0.881,0.119)$}{$(0,0,1)$}
{B}{参与归一化的指数为 $e^0=1$ 和 $e^2\approx7.389$。归一化后约为 $(0.119,0.881)$，掩码为 0 的位置保持为 0。}

% A05 S07:20 | 基础
\mc{主体行人附近有一位无关行人，但主体在较早时刻曾突然转向。若要分别建模这两类影响，HSGNN 的哪种分工最合适？}
{空间注意力选择重要行人，时间注意力选择重要历史轨迹}{空间注意力选择历史时刻，时间注意力选择道路节点}{空间注意力只处理群组，时间注意力只处理单个行人}{空间和时间注意力都只按行人间距离筛选}
{A}{空间交互图的自注意力选择影响主体行为的行人。时间交互图的自注意力选择影响主体未来行为的历史轨迹。}

% A06 S07:22 | 基础
\mc{两名行人的特征距离为 $M_{FS}^{1,2}=0.3$，阈值为 $\varepsilon=0.5$。按 $M_{1,2}=1/(1+e^{M_{FS}^{1,2}-\varepsilon})$ 计算，该权重最接近哪一个值？}
{$0.550$}{$0.450$}{$0.426$}{$0.622$}
{A}{$M_{1,2}=1/(1+e^{-0.2})\approx0.5498$，取三位小数为 $0.550$。}

% A07 S07:31 | 基础
\mc{静态群组的参数为 $d_{smin}=0.4$、$d_{xf}^L=1.6$、$d_{xb}^L=-1.2$、$d_{ym}^L=0.9$。按 $\sigma_{x1}^L=(d_{smin}+d_{xf}^L)/2$、$\sigma_{x2}^L=(d_{smin}+|d_{xb}^L|)/2$、$\sigma_y^L=(d_{smin}+d_{ym}^L)/2$ 计算 $(\sigma_{x1}^L,\sigma_{x2}^L,\sigma_y^L)$。}
{$(0.8,0.6,0.45)$}{$(1.0,1.0,0.65)$}{$(0.8,1.0,0.9)$}{$(1.0,0.8,0.65)$}
{D}{$\sigma_{x1}^L=(0.4+1.6)/2=1.0$，$\sigma_{x2}^L=(0.4+|-1.2|)/2=0.8$，$\sigma_y^L=(0.4+0.9)/2=0.65$。}

% A08 S07:30-31 | 基础
\mc{规划器要为静止行人与行走行人构造社交空间。为了匹配课件中的个体模型，哪项参数判断正确？}
{静态个体令 $v_h>0$，并令两个协方差相同}{动态个体令 $v_h=0$，并令 $\sigma_{x1}=\sigma_{x2}$}{静态个体令 $v_h=0$ 且边界连续，动态个体令 $v_h>0$ 且 $\sigma_{x1}>\sigma_{x2}$}{两类个体都只由行人数量 $N_h$ 决定}
{C}{课件以 $v_h=0$ 表示静态个体，并要求其两分量边界连续。动态个体满足 $v_h>0$，且前后方向尺度满足 $\sigma_{x1}>\sigma_{x2}$。}

% A09 S07:15 | 基础
\mc{实验报告称，加入 ACXCOOT 后计算耗时平均降低 $0.77$ s，即相对 COOT 降低 $46.10\%$。若按该比例反推，COOT 的平均耗时和新耗时最接近哪一组？}
{$1.67$ s 和 $0.90$ s}{$0.77$ s 和 $0.35$ s}{$1.67$ s 和 $1.67$ s}{$0.90$ s 和 $0.77$ s}
{A}{COOT 耗时约为 $0.77/0.4610=1.670$ s。加入 ACXCOOT 后约为 $1.670-0.77=0.900$ s。}

% A10 S07:24,27 | 基础
\mc{消融实验报告的 FDE 相对下降为：w/o GH $84.14\%$、w/o TIG $11.25\%$、w/o SS $39.94\%$。六人群组预测中，哪项消融结果提示群组层次对 FDE 的影响最大？}
{w/o TIG}{w/o SS}{w/o GH}{三者相同}
{C}{w/o GH 的 FDE 相对下降最大，为 $84.14\%$。这支持群组层次对该数据集的群组预测很重要，但不能单独证明失败原因。}

\exampagebreak
\sectiontitle{二、填空题}{10}{20}
\textbf{说明：}每题一个空，每空 2 分。数值保留 4 位小数时按题目要求作答。

% A11 S07:11 | 中等
\fillq{给定 $c_V=0.5$、$|x_g-x_c|=2.0$ m、$|y_g-y_c|=1.9$ m。按 $Tr_l=|x_g-x_c|/c_V$、$Tr_w=|y_g-y_c|/c_V$ 和 $h_A=[Tr_l-\lfloor Tr_w/2\rfloor+Tr_w]c_V$ 计算射束长度 $h_A=$\blankline。}
{$3.4$ m}{$Tr_l=4$，$Tr_w=3.8$，$\lfloor3.8/2\rfloor=1$。所以 $h_A=(4-1+3.8)\times0.5=3.4$ m。}

% A12 S07:10 | 中等
\fillq{ACXCOOT 取 $Iternum=b_2n_{pf}/(b_1e)$，并令 $Iter=\min(Iternum,Iter_{max})$。给定 $b_1=2$、$b_2=3$、$n_{pf}=40$、$Iter_{max}=20$，实际迭代次数为\blankline。}
{$20$}{未限幅值为 $120/(2e)\approx22.0728$。因其超过 $Iter_{max}=20$，实际值为 $20$。}

% A13 S07:22 | 中等
\fillq{当 $M_{FS}^{i,j}=0.3$、$\varepsilon=0.5$ 时，按 $M_{i,j}=1/(1+e^{M_{FS}^{i,j}-\varepsilon})$ 计算，$M_{i,j}$ 保留 4 位小数为\blankline。}
{$0.5498$}{代入得 $1/(1+e^{-0.2})=0.549834\ldots$，保留 4 位小数为 $0.5498$。}

% A14 S07:20 | 中等
\fillq{一行稀疏注意力的 logits 为 $(1,3,0)$，有效掩码为 $(1,1,0)$。对有效项减去最大 logit 后计算 softmax，无效项权重置零。输出权重保留 3 位小数为\blankline。}
{$(0.119,0.881,0)$}{有效项的指数为 $1$ 和 $e^2$。因此权重为 $(1/(1+e^2),e^2/(1+e^2),0)\approx(0.119,0.881,0)$。}

% A15 S07:21,23 | 中等
\fillq{一个行人的两个预测时刻真值为 $(0,0)$、$(2,0)$。某预测模式给出 $(0,1)$、$(2,1)$；另一模式给出 $(1,0)$、$(4,0)$。按 HSGNN 的 ADE 定义先对每个模式求平均距离，再取最小值，ADE 为\blankline。}
{$1$ m}{第一模式的距离和为 $1+1=2$ m，ADE 为 $2/2=1$ m。第二模式的距离和为 $1+2=3$ m，ADE 为 $1.5$ m。取较小值 $1$ m。}

% A16 S07:31 | 中等
\fillq{一个静态群组的三名行人在局部坐标中的 $x$ 坐标为 $0.6$、$-0.2$、$1.1$ m，$y$ 坐标为 $0.4$、$0.2$、$0.6$ m。按 $x_O^L=\sum_i x_i^L/N_h$、$y_O^L=\sum_i y_i^L/N_h$ 计算群组局部原点 $(x_O^L,y_O^L)=$\blankline。}
{$(0.5,0.4)$ m}{$x_O^L=(0.6-0.2+1.1)/3=0.5$ m，$y_O^L=(0.4+0.2+0.6)/3=0.4$ m。}

% A17 S07:33 | 中等
\fillq{轨迹速度为 $3$ m/s，$v_{max}=2$ m/s。按 $\delta_v=\|p'(t)\|^2-v_{max}^2$ 和罚项 $\max(\delta_v,0)^3$ 计算该时刻速度罚项为\blankline。}
{$125$}{$\delta_v=3^2-2^2=5$。因此罚项为 $\max(5,0)^3=125$。}

% A18 S07:33 | 中等
\fillq{某局部轨迹的平滑项为 $J_{ct}=12$，总时间为 $J_t=20$ s，时间权重为 $\rho=0.25$。按 $J=J_{ct}+\rho J_t$ 计算总目标值为\blankline。}
{$17$}{代入得 $J=12+0.25\times20=17$。本题只计算题面给定的两个目标项。}

% A19 S07:35 | 中等
\fillq{某个社交空间的 $\phi_{\mu,\hat\Sigma}(p_{bot})/A=0.4$，$\theta_{hr}=\pi/2$，$\theta_{rh}=-\pi/2$。按 $SGI=[\phi_{\mu,\hat\Sigma}(p_{bot})/A][1-0.9\theta_{hr}/\pi][1-\theta_{rh}/\pi]$ 计算该空间的 SGI 项为\blankline。}
{$0.33$}{$0.4[1-0.9(1/2)][1-(-1/2)]=0.4\times0.55\times1.5=0.33$。}

% A20 S07:24 | 中等
\fillq{消融实验中 w/o GH、w/o TIG、w/o SS 的 FDE 相对下降分别为 $84.14\%$、$11.25\%$、$39.94\%$。若按 FDE 变化幅度排序，影响最大的是移除\blankline。}
{群组划分模块（GH）}{$84.14\%$ 大于 $39.94\%$ 和 $11.25\%$，因此影响最大的是移除 GH。此比较描述该组消融结果，不表示所有数据集都具有相同排序。}

\exampagebreak
\sectiontitle{三、问答题}{3}{30}

% A21 S07:10-11,15 | 中等
\essay{ETPP 在道路图上使用带信息增益项的 Dijkstra 变形，并使用 ACXCOOT 搜索边界点遍历顺序。按下列给定数值完成一次递进计算。
\begin{subquestions}
\item 当前节点的累计代价 $g_V=4$、$c_V=2$、$D_{free}(V_c)=0$、$\chi_1=0.8$。邻居 A 的距离为 $D_A=3$、$D_{free}(A)=2$、$\chi_{2,A}=0.6$；邻居 B 的距离为 $D_B=2$、$D_{free}(B)=0$、$\chi_{2,B}=0.5$。分别计算 $H_{IE}$ 和扩展代价 $f_V$，并按较小代价选择下一节点。（4 分）
\item 若边界点概率为 $(0.5,0.5,0,0)$，地图最大熵为 $\ln4$，停止阈值为 $I_{thr}=0.6$，计算归一化熵并判断是否停止探索。（3 分）
\item 取 $b_1=2$、$b_2=3$、$n_{pf}=40$、$Iter_{max}=20$。计算限幅前后的 ACXCOOT 迭代次数，并说明自适应项随边界点数如何变化。（3 分）
\end{subquestions}}
{\textbf{参考答案与评分：}
\begin{subquestions}
\item A、B 的 $H_{IE}$ 分别为 $\min(0.8,0.6e^{-2/2})=0.2207$ 和 $\min(0.8,0.5)=0.5$，各得 0.5 分。$f_A=4+(1+0.2207)3=7.6621$、$f_B=4+(1+0.5)2=7$，各得 1 分。比较代价并选 B，得 1 分。
\item $H=-\sum p_i\ln p_i=\ln2$，得 1 分。$I=\ln2/\ln4=0.5$，得 1 分。$0.5\le0.6$，按题面判据停止探索，得 1 分。
\item 限幅前 $Iternum=3\times40/(2e)\approx22.0728$，得 1 分。限幅后为 $\min(22.0728,20)=20$，得 1 分。其他参数固定时，$n_{pf}$ 增大将使未限幅迭代次数按比例增大，得 1 分。
\end{subquestions}}

% A22 S07:19-27 | 综合
\essay{HSGNN 用群组划分和稀疏时空图学习预测行人轨迹。某帧有 4 名行人，特征距离矩阵上三角元素为 $M_{FS}^{1,2}=0.3$、$M_{FS}^{1,3}=0.8$、$M_{FS}^{1,4}=1.1$、$M_{FS}^{2,3}=0.4$、$M_{FS}^{2,4}=0.9$、$M_{FS}^{3,4}=0.2$，分组阈值 $\varepsilon=0.5$。
\begin{subquestions}
\item 按 $M_{FS}^{i,j}\le\varepsilon$ 列出群组划分模块保留的无序行人对。（2 分）
\item 按 $M_{i,j}=1/(1+e^{M_{FS}^{i,j}-\varepsilon})$ 计算 $(1,2)$ 与 $(3,4)$ 两对的权重，保留 3 位小数。（2 分）
\item 主体旁有一名关键行人，主体在较早时刻有一次突然变向。分别指出空间和时间交互图应关注的对象，并说明为什么时间分支不做时空融合。（2 分）
\item w/o GH 与 w/o SS 的 ADE/FDE 相对下降分别为 $65.81\%/84.14\%$ 和 $44.25\%/39.94\%$。比较 FDE 变化，并为六人群组划分失败提出两种可检验原因及一项后续实验。（4 分）
\end{subquestions}}
{\textbf{参考答案与评分：}
\begin{subquestions}
\item 保留 $(1,2)$、$(2,3)$、$(3,4)$ 三对。每对正确得 $2/3$ 分，共 2 分。题面按阈值直接判断成对关系，不要求推断传递闭包。
\item $M_{1,2}=1/(1+e^{-0.2})\approx0.550$，得 1 分。$M_{3,4}=1/(1+e^{-0.3})\approx0.574$，得 1 分。
\item 空间分支关注影响主体行为的关键行人，得 1 分。时间分支关注影响未来行为的历史轨迹；行人数量在每次预测中不固定，因此不做时空融合，得 1 分。
\item w/o GH 的 FDE 下降 $84.14\%$，大于 w/o SS 的 $39.94\%$，得 1 分。原因可为大群组样本不足或网络未学到群组深层特征，说明两项各得 1 分。后续实验可补充不同规模群组数据，或在分层控制组上比较分组质量与轨迹误差，提出可验证方案得 1 分。
\end{subquestions}}

% A23 S07:29-33 | 综合
\essay{某移动机器人要经过一名静止行人附近。采用单个静态高斯社交空间，令 $\mu=(0,0)$、$\Sigma=\operatorname{diag}(1,1)$。用相对场 $r(p)=\exp[-\frac12(p-\mu)^T\Sigma^{-1}(p-\mu)]$ 表示空间强度，并以 $r(p)>e^{-1}$ 判为进入社交空间。为比较候选轨迹，题面定义离散平滑项 $J_{ct}=\sum_{t=1}^{T-2}\|p_{t+2}-2p_{t+1}+p_t\|^2$ 和目标 $J=J_{ct}+\rho T$，取 $\rho=1$。
\begin{subquestions}
\item 候选 A 的采样点为 $(0,2),(0,1),(0,2)$。计算中间点的 $r(p)$，判断候选 A 是否满足社交空间约束。（3 分）
\item 候选 B 为 $(2,0),(2,1),(2,2)$，候选 C 为 $(3,0),(2,0),(3,0)$。检查两条轨迹的采样点约束，并计算离散平滑项。（3 分）
\item 计算 B、C 的目标值并选择较优候选。结合课件框架，说明为何还需检查采样点之间的连续轨迹、障碍物和运动学约束。（4 分）
\end{subquestions}}
{\textbf{参考答案与评分：}
\begin{subquestions}
\item 中间点 $(0,1)$ 的马氏距离平方为 1，所以 $r=e^{-1/2}\approx0.6065$，得 2 分。$0.6065>e^{-1}\approx0.3679$，候选 A 进入社交空间，不满足约束，得 1 分。
\item B 的三个点到原点的距离平方分别为 4、5、8，故 $r$ 均小于 $e^{-1}$。其二阶差分为 $(0,0)$，所以 $J_{ct,B}=0$，得 1.5 分。C 的距离平方为 9、4、9，也满足采样点约束。其二阶差分为 $(2,0)$，所以 $J_{ct,C}=4$，得 1.5 分。
\item 两条轨迹均有 $T=3$。$J_B=0+3=3$，$J_C=4+3=7$，应选 B，计算和选择共 2 分。有限采样可能漏掉采样点之间的碰撞或社交空间穿越，得 1 分。课件目标还包含障碍物避碰及速度、加速度、角速度等运动学约束，说明仍需检查其中至少两项得 1 分。
\end{subquestions}}

\exampagebreak
\sectiontitle{四、程序题}{2}{30}

% A24 S07:20 | 中等
\programq{实现稀疏行 softmax 函数 \code{sparse_softmax(logits, active)}，用于生成 HSGNN 的稀疏交互权重。

\textbf{输入：}\code{logits} 是有限浮点数 NumPy 数组，形状为 $(R,C)$，其中 $R,C\ge1$。\code{active} 是同形状的布尔 NumPy 数组。输入满足要求。

\textbf{输出与规则：}
\begin{subquestions}
\item 返回同形状浮点数组。对每一行，仅对 \code{active} 为 \code{True} 的元素计算 softmax；其余位置必须为 0。（5 分）
\item 每行至少有一个有效项时，有效权重和为 1。请使用减去该行有效 logits 最大值的方式计算，以避免指数溢出。（5 分）
\item 若某行没有有效项，该行全为 0。题面自检输入应产生第一行约为 $(0.1192,0.8808,0)$、第二行全零。（5 分）
\end{subquestions}
\ifanswers\else
\textbf{可用 API：}
\codefile{snippets/A01.py}
\fi
输入数组形状有效，无需增加其他依赖。}
{\textbf{参考程序：}
\codefile{programs/A24.py}
\textbf{评分：}输出形状及无效项置零得 3 分；有效项按行稳定归一化得 5 分；空行全零得 3 分；样例结果和权重和自检正确得 4 分。}

% A25 S07:30-33 | 综合
\programq{实现社会空间约束下的候选路径选择函数 \code{select_social_path(paths, durations, means, sigmas, weights, threshold, rho)}。

\textbf{输入：}\code{paths} 是有限浮点数组，形状为 $(K,T,2)$，表示 $K$ 条候选路径、每条路径有 $T\ge3$ 个二维采样点；$K$ 可以为 0。\code{durations} 长度为 $K$，每个值为正。\code{means} 和 \code{sigmas} 形状均为 $(G,2)$，其中 $G\ge1$，每个标准差为正。\code{weights} 长度为 $G$，非负且和为 1。\code{threshold} 为正，\code{rho} 非负。所有输入均有限且形状有效。

每个社交空间高斯分量采用对角协方差。位置 $p$ 的混合密度为
\[
\phi(p)=\sum_{g=1}^{G}\frac{w_g}{2\pi\sigma_{gx}\sigma_{gy}}
\exp\!\left[-\frac12\sum_{d\in\{x,y\}}\left(\frac{p_d-\mu_{gd}}{\sigma_{gd}}\right)^2\right].
\]

\textbf{规则：}
\begin{subquestions}
\item 若某路径任一采样点满足 $\phi(p)>\mathrm{threshold}$，则该路径进入社交空间并应排除。等于阈值时保留。（5 分）
\item 对保留路径计算 $J_{ct}=\sum_{t=0}^{T-3}\|p_{t+2}-2p_{t+1}+p_t\|^2$，并计算 $J=J_{ct}+\rho\,\mathrm{duration}$。（5 分）
\item 返回代价最小路径的二元组（原始路径下标，$J$）。代价相同时选下标较小者。没有可行路径时返回 \code{None}。（5 分）
\end{subquestions}
\ifanswers\else
\textbf{可用 API：}
\codefile{snippets/A02.py}
\fi
本题把连续轨迹的平滑积分离散为采样点二阶差分平方和。无需实现完整轨迹优化器。}
{\textbf{参考程序：}
\codefile{programs/A25.py}

\textbf{自检结果：}第一条路径在 $(0,0)$ 处的混合密度大于阈值，因此排除。第二条路径为直线，$J_{ct}=0$，持续时间为 4，故 $J=4$。第三条路径的二阶差分为 $(2,0)$，故 $J_{ct}=4$；持续时间为 2，故 $J=6$。返回 \code{(1, 4.0)}。

\textbf{评分：}混合高斯密度及所有采样点检查正确得 5 分；二阶差分与时间项正确得 5 分；最小值、平局规则和无可行路径处理正确得 3 分；自检下标与代价正确得 2 分。}
